\section{第07章 ソフトウェア保守}
\begin{pointbox}{この資料の主張}ソフトウェア保守は、単なる「バグ修正」ではない。運用中のソフトウェアを価値ある状態に保ち、環境変化、利用者要求、障害、セキュリティ脅威、技術的陳腐化に対応し続ける活動である。保守は納品後だけに始まるものではなく、開発初期から保守性、移行、運用、支援、測定を計画する必要がある。\end{pointbox}
\begin{pointbox}{この資料の読み方}専門用語は原則として日本語で示し、必要に応じて英語を括弧内に併記する。英語名は、元資料やツール名を確認する手がかりとして残す。本文は初学者が前提知識なしに読めるように、概念の目的、実務で見るべき対象、用語の境界を重視している。\end{pointbox}
この資料は SWEBOK Guide v4.0a Chapter 07 の内容を学習用に要約・再構成したものであり、逐語訳ではない。\par
